%% ma: L12-B08-0175
%% lop: 12
%% chuong: 2
%% bai: 8
%% dang: 12.8.30
%% loai: TN4
%% muc: VD
%% dap_an: "B"
%% nguon: "CĐ2.4.pdf giáo viên gửi 27/10/2026 (M07, Lê Minh Khôi – Nguyễn Anh Khôi, Chương II Chủ đề 4 Các bài toán về tâm tỉ cự), Câu 25"
%% kien_thuc: "Bài 8 (khoảng cách hai điểm, biểu thức tọa độ), Bài 6 (bình phương vô hướng, tâm tỉ cự)"
%% trang_thai: chinh-thuc
\begin{ex}
Trong không gian $Oxyz$, cho các điểm $A(2;4;-1)$, $B(1;4;-1)$, $C(2;4;3)$, $D(2;2;-1)$. Gọi $M(x;y;z)$ là điểm thỏa mãn $MA^2+MB^2+MC^2+MD^2$ đạt giá trị nhỏ nhất. Tính giá trị của biểu thức $P=x+y+z$.
\choice{$P=6$}{\True $P=\dfrac{21}4$}{$P=8$}{$P=9$}
\loigiai{Tổng nhỏ nhất khi $M$ là trọng tâm $G$ của tứ diện: $G\left(\dfrac74;\dfrac72;0\right)$. Vậy $P=\dfrac{21}4$.}
\end{ex}
